\section{Exact section shear, quotient section and the full
receiver}\label{exact-section-shear-quotient-section-and-the-full-receiver}

25 September 2026. Independent derivation SSR0--SSR7. This note reads
the complete current NPE0--10 and NER0--11 sources and verifies the
proposed replacement of the residue section. It does not edit either
source. Its conclusions concern the actual coefficient extension and do
not identify it with Deligne's geometric lifting sequence.

\subsection{SSR0. Source, construction stage and reading
coverage}\label{ssr0.-source-construction-stage-and-reading-coverage}

The supporting datum remains
\(\tau\langle Z_1;\mathrm{no}\ Z_2\rangle\). All coefficient operations
below follow the complete-history arithmetic reconstruction. They assign
no arithmetic value, coordinate, parity, metric or retracted addition to
that supporting datum. Separate branch histories remain separate.

Use exactly the original Fréchet source \[
\mathcal B=\{F\in\mathcal O(\mathbb C):q_{A,N}(F)<\infty\},\qquad
q_{A,N}(F)=\sup_{|\Re s|\le A}(1+|\Im s|)^N|F(s)|,
\tag{SSR0.1}
\] with all \(A>0\) and integers \(N\ge0\). Its closed ideal \(I\)
imposes the complete original nontrivial-zero jets through order
\(m_\rho-1\); \(Q=\mathcal B/I\). Keep \[
F_0(s)=\frac{s(s-1)}8\pi^{-s/2}\Gamma(s/2)\zeta(s),\qquad
F_0(0)=\frac18,\qquad g_t(s)=e^{ts^2},\quad t>0.
\tag{SSR0.2}
\] The equality retains all original factors and exceptional-point
contributions in NPE0.2; it does not replace original \(\zeta\). The
current comparison uses the already established membership
\(F_0\in\mathcal B\), the full original-zero ideal, and NPE1's
meromorphic extension \[
h=h_t=8g_tF_0/s,\qquad E=\mathcal B\oplus\mathbb C h,
\qquad\operatorname{res}_0h=1.
\tag{SSR0.3}
\] Here and throughout, a product decomposition of \(E\) denotes
coordinates on actual meromorphic functions with the transported product
topology. It does not assert closure of \(E\) under pointwise
multiplication.

Complete source reading in this independent task:

\begin{itemize}
\tightlist
\item
  \nolinkurl{NOOR_MEROMORPHIC_ENDPOINT_EXTENSION.md}, NPE0--10, SHA256
  \nolinkurl{da8a4d25c22391e32fddc44fac98d12de515679f2d5b013ee33639c6200134ed}.
\item
  \nolinkurl{NOOR_ENDPOINT_EXTENSION_INDEPENDENT_REVIEW.md}, NER0--11,
  SHA256
  \nolinkurl{c05f47854508143b2aeb3ea42a7e0dbd74fe5c3da9b00923dd0bacc6ad44442e}.
\end{itemize}

The calculations below are direct comparisons of their exact maps. No
additional human-source reading is claimed.

\subsection{SSR1. The two residue sections give a continuous
shear}\label{ssr1.-the-two-residue-sections-give-a-continuous-shear}

Define \[
j=j_t=g_t/s,\qquad
k=k_t=j-h=\frac{g_t(1-8F_0)}s.
\tag{SSR1.1}
\] The numerator of \(k\) is an element of \(\mathcal B\): both \(g_t\)
and \(8g_tF_0\) are. It vanishes at zero by (SSR0.2), so division by
\(s\) is entire. For any \(K\in\mathcal B\) with \(K(0)=0\), the maximum
principle on the radius-two disk, together with division outside the
unit disk, gives \[
q_{A,N}(K/s)\le q_{A,N}(K)+2^{N-1}q_{2,0}(K).
\tag{SSR1.2}
\] Indeed outside the unit disk \(|1/s|\le1\); inside it,
\(|K(s)/s|\le\tfrac12\sup_{|z|=2}|K(z)|\), and the vertical weight is at
most \(2^N\). Thus \(k\in\mathcal B\), and \(j=h+k\in E\) has residue
one. Consequently \[
E=\mathcal B\oplus\mathbb C j
\tag{SSR1.3}
\] as the same meromorphic space, with the same topology. In fact, if \[
f=F+ch=G+cj,
\] then the two coordinate maps are exactly \[
(F,c)_h\longmapsto(G,c)_j=(F-ck,c),\qquad
(G,c)_j\longmapsto(G+ck,c)_h.
\tag{SSR1.4}
\] Both maps are continuous linear inverses because \(k\) is a fixed
member of \(\mathcal B\).

The endpoint coefficients verify the sign independently: \[
j(s)=s^{-1}+ts+O(s^3),\qquad
h(s)=s^{-1}+8F_0'(0)+O(s),\qquad k(0)=-8F_0'(0).
\tag{SSR1.5}
\] Thus the Laurent constant of \(f\) is \(F(0)+8cF_0'(0)=G(0)\). The
endpoint datum has been retained under the shear.

\subsection{SSR2. Cover cocycles and strong-dual
coordinates}\label{ssr2.-cover-cocycles-and-strong-dual-coordinates}

For every recovered integer \(n\ge1\), let \(U_nf=n^{1-s}f\). Retain the
old cocycle and define the new one: \[
\delta_n=nh-U_nh=8g_tF_0\frac{n-n^{1-s}}s\in I,
\qquad
\eta_n=nj-U_nj=g_t\frac{n-n^{1-s}}s\in\mathcal B.
\tag{SSR2.1}
\] Membership follows from the same removable division estimate as
(SSR1.2); the numerator is in \(\mathcal B\) and vanishes at zero. Both
have value \(n\log n\) at zero. The exact comparison, including its
sign, is \[
\boxed{\eta_n=\delta_n+(n-U_n)k.}
\tag{SSR2.2}
\] It follows by substituting \(j=h+k\) into \(nj-U_nj\). In the
respective coordinates the same actual operator is \[
U_n^E(F,c)_h=(U_nF-c\delta_n,nc)_h,
\qquad
U_n^E(G,c)_j=(U_nG-c\eta_n,nc)_j.
\tag{SSR2.3}
\] Both cocycles satisfy \[
c_{mn}=m c_n+U_n c_m=n c_m+U_m c_n
\quad(c=\delta\text{ or }\eta).
\tag{SSR2.4}
\] For example \(m(nj-U_nj)+U_n(mj-U_mj)=mnj-U_{mn}j\). This proves the
formula without discarding either summand.

Write an actual continuous complex-linear functional on \(E\) as \[
\widetilde\Lambda(F+ch)=\Lambda(F)+c\alpha,
\qquad
\widetilde\Lambda(G+cj)=\Lambda(G)+c\beta.
\tag{SSR2.5}
\] The restriction \(\Lambda\) to the original entire subspace is
unchanged. Since \(\widetilde\Lambda(j)=\widetilde\Lambda(h+k)\), the
exact dual shear is \[
\boxed{\beta=\alpha+\Lambda(k),\qquad\alpha=\beta-\Lambda(k).}
\tag{SSR2.6}
\] It is a strong-dual topological isomorphism. Evaluation at a fixed
\(k\) is continuous for the strong topology because \(\{k\}\) is
bounded, and the displayed inverse is continuous for the same reason.
The underlying product identifications follow because product-bounded
sets have bounded projections and bounded factor sets embed in the
product.

Transposing (SSR2.3), with no conjugation in the complex-linear
transpose, gives \[
(U_n^E)'(\Lambda,\alpha)_h
  =(U_n'\Lambda,n\alpha-\Lambda(\delta_n))_h,
\] \[
\boxed{(U_n^E)'(\Lambda,\beta)_j
  =(U_n'\Lambda,n\beta-\Lambda(\eta_n))_j.}
\tag{SSR2.7}
\] Substituting (SSR2.6) and then (SSR2.2) gives the second formula from
the first: the new final coordinate is
\(n\alpha-\Lambda(\delta_n)+\Lambda(U_nk)=n\beta-\Lambda(\eta_n)\).

\subsection{SSR3. The original full-ideal quotient keeps its unique
equivariant
section}\label{ssr3.-the-original-full-ideal-quotient-keeps-its-unique-equivariant-section}

Because \(I\) is closed in \(\mathcal B\), either product coordinate
system gives a topological isomorphism \(E/I\simeq Q\oplus\mathbb C\).
The change is the induced shear \[
([F],c)_h\longmapsto([F]-c[k],c)_j.
\tag{SSR3.1}
\] In old coordinates the quotient action is diagonal, because
\(\delta_n\in I\). In new coordinates it is \[
([G],c)\longmapsto([U_nG]-c[\eta_n],nc),
\qquad [\eta_n]=(n-U_n)[k].
\tag{SSR3.2}
\] The equivariant residue section furnished by the old coordinates is
therefore \[
c\longmapsto c[h]=c[j-k],
\qquad c\longmapsto(-c[k],c)_j,
\tag{SSR3.3}
\] where the first expression is understood as \(c[h]\). Direct
substitution in (SSR3.2) gives \((-nc[k],nc)\), proving equivariance.

For completeness, this section is unique already for each recovered
\(n>1\). Any residue section is \(c\mapsto c([j]+[a])\) for some
\([a]\in Q\). Its equivariance is precisely \[
(U_n-n)[a]=[\eta_n],
\quad\text{equivalently}\quad
(U_n-n)([a]+[k])=0.
\tag{SSR3.4}
\] At every actual nontrivial zero \(\rho\), \(0<\Re\rho<1\) gives
\(|n^{1-\rho}|=n^{1-\Re\rho}<n\). Hence \(n^{1-s}-n\) is a holomorphic
unit near that point. If its product with a representative \(a+k\) lies
in \(I\), division by this unit preserves all vanishing derivatives
through order \(m_\rho-1\), so \(a+k\in I\). Thus \([a]=-[k]\), proving
uniqueness with all multiplicities retained.

The new section \([j]\) itself is not equivariant: for every actual
nontrivial zero, \[
\eta_n(\rho)=\frac{g_t(\rho)}{\rho}
                  (n-n^{1-\rho})\ne0\quad(n>1).
\tag{SSR3.5}
\] The same inequality and \(\rho\ne0\), \(g_t(\rho)\ne0\) prove this
exact assertion. In particular \(\eta_n\notin I\) and \([k]\ne0\) in the
actual zeta quotient. This uses the actual nonempty nontrivial divisor
already present in the source, not a conjectured off-line zero.

There is no equivariant residue section in \(E\) itself, whether written
using \(h\) or \(j\): its new-coordinate equation would be
\((U_n-n)a=\eta_n\), whose left side vanishes at zero whereas the right
side has value \(n\log n\ne0\). The unique section in (SSR3.3) is
specifically a section in \(E/I\).

\subsection{SSR4. Full jets and the old zero-coordinate lift after the
shear}\label{ssr4.-full-jets-and-the-old-zero-coordinate-lift-after-the-shear}

For an original nontrivial zero \(\rho\), and \(0\le d<m_\rho\), the
complete zero order of \(h\) gives \[
k^{(d)}(\rho)=j^{(d)}(\rho)
=\sum_{a=0}^d\binom da g_t^{(a)}(\rho)
       (-1)^{d-a}(d-a)!\rho^{-(d-a+1)}.
\tag{SSR4.1}
\] This is the full Leibniz formula; no derivative is replaced by its
value component. Consequently the actual derivative functional
\(\operatorname{ev}_\rho^{(d)}\) on \(E\) has coordinates \[
(\operatorname{ev}_\rho^{(d)}|_{\mathcal B},0)_h
\longleftrightarrow
(\operatorname{ev}_\rho^{(d)}|_{\mathcal B},k^{(d)}(\rho))_j.
\tag{SSR4.2}
\] In particular the zeroth new coordinate is \(g_t(\rho)/\rho\ne0\),
even though the old coordinate is zero.

More generally, the old lift \(\sigma_h\Lambda=(\Lambda,0)_h\) becomes
\[
\sigma_h\Lambda=(\Lambda,\Lambda(k))_j.
\tag{SSR4.3}
\] It is the graph of evaluation on \(k\), not the new zero-coordinate
slice. For \(\Lambda\in I^\perp\), (SSR2.2) gives \[
n\Lambda(k)-\Lambda(\eta_n)
=\Lambda(U_nk)=(U_n'\Lambda)(k),
\tag{SSR4.4}
\] because \(\Lambda(\delta_n)=0\). Thus this tilted graph is invariant
on the full original quotient dual. Conversely, for a vector on that
graph its transformed failure to stay on the graph is exactly
\(-\Lambda(\delta_n)\), directly by (SSR2.2). This preserves the old
complete ideal criterion; no statement about the distinct new
zero-coordinate slice follows merely by renaming coordinates.

\subsection{SSR5. The full uncorrected receiver in the new
coordinates}\label{ssr5.-the-full-uncorrected-receiver-in-the-new-coordinates}

Retain the original meromorphic tests, including the exceptional zeroth
coefficient: \[
\phi_0(s)=-1/s,\qquad
\phi_m(s)=\frac{m^{1-s}-(m+1)^{1-s}}s\ (m\ge1),
\qquad G_m=g_t\phi_m.
\tag{SSR5.1}
\] Each has residue \(-1\). Define the two entire representatives \[
F_m^h=G_m+h,\qquad F_m^j=G_m+j=F_m^h+k.
\tag{SSR5.2}
\] They belong to \(\mathcal B\) by NPE5 and SSR1, with the same
polynomial-in-\(m\) seminorm estimates up to a fixed added term. The
zeroth term changes explicitly: \[
F_0^h=-k,\qquad F_0^j=0,
\] \[
F_m^j=\frac{g_t(s)}s
\bigl(m^{1-s}-(m+1)^{1-s}+1\bigr)\quad(m\ge1).
\tag{SSR5.3}
\] For \(\Lambda\in\mathcal B'\), the corresponding disk-holomorphic
receivers obey \[
H_j\Lambda(z)=\sum_{m\ge0}\overline{\Lambda(F_m^j)}z^m
=H_h\Lambda(z)+\frac{\overline{\Lambda(k)}}{1-z}.
\tag{SSR5.4}
\] These are conjugate-linear receiver maps; this conjugation is
distinct from the unconjugated transpose in SSR2. The geometric series
converges on every compact subdisk. The full uncorrected receiver is
therefore exactly \[
\boxed{M_t\widetilde\Lambda
=H_h\Lambda-\frac{\overline\alpha}{1-z}
=H_j\Lambda-\frac{\overline\beta}{1-z}.}
\tag{SSR5.5}
\] The last equality uses \(\beta=\alpha+\Lambda(k)\). The new tests'
zeroth value gives the useful exact identity \[
H_j\Lambda(0)=0,\qquad
M_t\widetilde\Lambda(0)=-\overline\beta,
\qquad\beta=-\overline{M_t\widetilde\Lambda(0)}.
\tag{SSR5.6}
\] Thus the new scalar coordinate is recovered from the full receiver's
constant coefficient. In particular, the added boundary term has not
disappeared.

For the original quotient dual lift, substitute (SSR4.3) into (SSR5.5):
\[
M_t\sigma_h\Lambda
=H_j\Lambda-\frac{\overline{\Lambda(k)}}{1-z}
=H_h\Lambda\qquad(\Lambda\in I^\perp).
\tag{SSR5.7}
\] All old received functions are retained exactly. Setting \(\beta=0\)
instead would omit the generally nonzero term shown in (SSR5.7).

\subsection{SSR6. Cover identity and maximal Hardy domain under the
shear}\label{ssr6.-cover-identity-and-maximal-hardy-domain-under-the-shear}

Telescoping the complete tests in (SSR5.1), including \(m=0\), gives \[
\sum_{a=0}^{n-1}G_{nm+a}=U_nG_m,
\quad
\sum_{a=0}^{n-1}F_{nm+a}^j=U_nF_m^j+\eta_n.
\tag{SSR6.1}
\] The second follows by adding \(nj\) and using \(nj-U_nj=\eta_n\); its
first block also gives \[
\sum_{a=0}^{n-1}F_a^j=\eta_n
\tag{SSR6.2}
\] because \(F_0^j=0\). Let \(\mathscr W_n^*\) denote the original
coefficient-block sum on disk-holomorphic functions. Taking the
functionals with the conjugations of (SSR5.4) proves \[
\mathscr W_n^*H_j\Lambda-H_jU_n'\Lambda
=\frac{\overline{\Lambda(\eta_n)}}{1-z}.
\tag{SSR6.3}
\] Since \(\mathscr W_n^*(1-z)^{-1}=n(1-z)^{-1}\), the complete (SSR5.5)
and (SSR2.7) give \[
\mathscr W_n^*M_t=M_t(U_n^E)'
\tag{SSR6.4}
\] in the new coordinates too. Both terms in the last expression of
(SSR5.5) are retained in this exact covariance calculation.

The same maximal Hardy graph is transported by the shear as \[
\{(\Lambda,\alpha):H_h\Lambda-\overline\alpha/(1-z)\in H^2\}
\longleftrightarrow
\{(\Lambda,\beta):H_j\Lambda-\overline\beta/(1-z)\in H^2\}.
\tag{SSR6.5}
\] This is equality of actual received functions under the continuous
dual shear, so it is a homeomorphism for the graph topologies as well.
It changes neither injectivity nor the retained raw-cover degree
identity in NPE7.3. In particular, it gives no permission to assert
Hardy membership for \(H_j\Lambda\) alone when only the difference in
(SSR6.5) is known to belong to \(H^2\).

\subsection{SSR7. Review conclusion and exact extension of the
calculation}\label{ssr7.-review-conclusion-and-exact-extension-of-the-calculation}

All shear formulas hold more generally for any two continuous residue
sections differing by a fixed \(k\in\mathcal B\): the primal shear is
subtraction of \(ck\), the dual shear is addition of \(\Lambda(k)\), and
the cocycles differ by \((n-U_n)k\). Their proofs are exactly the
substitutions in SSR1--2 and do not require Gaussian or zeta factors.
For the specific sections examined here, SSR3--6 additionally retain the
complete original zero jets and full uncorrected receiver. These extra
conclusions use the actual original ideal and cannot be inferred for an
arbitrary quotient.

No mathematical error was found in the supplied proposed section change
once the two coordinate systems are kept distinct. The sign of the dual
shear is positive; the cover transpose contains the negative cocycle
evaluation. The equivariant residue section is unique in the full-ideal
pushout and is \([h]=[j-k]\); it is not \([j]\) and it is not an
equivariant section in \(E\). The existing original quotient-dual lift
is the graph \(\beta=\Lambda(k)\) in the new coordinates. These are
exact comparisons of the same source and action, not a deletion of its
endpoint or original-zeta data.
